% Article: The parameter-derivative tower of the generalized Stieltjes constants
% - Created (UTC): 2026-06-28T23:08:52Z
% - Repository HEAD: f594924c0dcc98e0fcb8147b287a09aa52f75955
% Source corpus: src/PolyLog/docs/reports/
%   stieltjes-general-k-duality-and-gamma2-tables__a4f0bee12e5b.md,
%   stieltjes-second-parameter-derivative-s3-layer__221c04334769.md,
%   stieltjes-derivative-relations__41afbe98f668.md, and the supporting tools
%   src/PolyLog/tools/stieltjes_general_k.py (+ scratch/g1k_grid_rank.py,
%   stieltjes_s2deriv_hunts.py, stieltjes_s3.py).
\documentclass[11pt]{article}
\usepackage[margin=1.1in]{geometry}
\usepackage{amsmath,amssymb,amsthm}
\usepackage{booktabs}
\usepackage{microtype}
\usepackage{xcolor}
\usepackage[colorlinks=true,linkcolor=blue!50!black,citecolor=blue!50!black,urlcolor=blue!60!black]{hyperref}
\usepackage{xurl}
\hypersetup{
  pdftitle={The Parameter-Derivative Tower of the Generalized Stieltjes
    Constants: a Uniform Harmonic-Number Theorem},
  pdfauthor={Vladimir Reshetnikov},
}

\newtheorem{theorem}{Theorem}[section]
\newtheorem{proposition}[theorem]{Proposition}
\newtheorem{lemma}[theorem]{Lemma}
\newtheorem{corollary}[theorem]{Corollary}
\theoremstyle{remark}
\newtheorem{remark}[theorem]{Remark}

\newcommand{\Glaisher}{\mathsf{A}}
\newcommand{\Cat}{\mathsf{G}}
\newcommand{\Z}{\mathbb{Z}}
\newcommand{\Q}{\mathbb{Q}}
\newcommand{\h}[1]{H_{#1}}
\newcommand{\zd}[2]{\zeta^{(#1)}(#2)}      % j-th s-derivative of zeta
\newcommand{\Ld}[2]{L^{(#1)}(#2)}
\newcommand{\chim}[1]{\chi_{-#1}}
\DeclareMathOperator{\Li}{Li}

\title{The Parameter-Derivative Tower of the Generalized\\
Stieltjes Constants: a Uniform Harmonic-Number Theorem}
\author{Vladimir Reshetnikov\thanks{%
v.reshetnikov@gmail.com. Prepared with agentic assistance (Anthropic Claude) in
the \emph{Smithereens} repository,
\texttt{github.com/VladimirReshetnikov/Smithereens}, \texttt{src/PolyLog}.
Provenance: created 2026-06-28 (UTC); repository snapshot
\texttt{f594924c0dcc98e0fcb8147b287a09aa52f75955}. Every displayed identity was
verified numerically (mpmath, \(60\)--\(220\) digits), the \(q=4\) entries
independently in a second computer-algebra system (Wolfram, \(\sim426\)
digits), and the whole development was re-derived by an independent multi-agent
workflow.}}
\date{June 2026}

\begin{document}
\maketitle

\begin{abstract}
The generalized Stieltjes constants $\gamma_n(a)$, regarded as functions of the
Hurwitz parameter $a$, have a parameter-derivative tower whose $k$-th rung lives
at the $s=k+1$ derivative layer of the Hurwitz zeta function. We show this tower
is \emph{uniform in $k$} and governed entirely by the harmonic numbers: writing
$\h{k}=\sum_{m\le k}1/m$ and $\h{k}^{(2)}=\sum_{m\le k}1/m^2$,
\[
  \gamma_1^{(k)}(a)=(-1)^{k+1}k!\,\bigl[\zeta'(k{+}1,a)+\h{k}\,\zeta(k{+}1,a)\bigr],
\]
\[
  \gamma_2^{(k)}(a)=(-1)^{k}k!\,\bigl[\zeta''(k{+}1,a)+2\h{k}\,\zeta'(k{+}1,a)
   +(\h{k}^2-\h{k}^{(2)})\,\zeta(k{+}1,a)\bigr],
\]
and likewise for every $n$, with the elementary symmetric functions of
$\{1,\tfrac12,\dots,\tfrac1k\}$ as coefficients. The same $\h{k}$ controls the
rational-argument landscape: a character coordinate
$\sum_p\chi(p)\gamma_1^{(k)}(p/q)=(-1)^{k+1}k!\,q^{k+1}[L'(k{+}1,\chi)+(\h{k}+\ln
q)L(k{+}1,\chi)]$, a distribution row, and a residual-lattice dimension
$\varphi(q)-1$ that is \emph{independent of $k$}. We give explicit closed-form
tables for the second-constant derivatives $\gamma_2'(p/q)$ and
$\gamma_2''(p/q)$ at $q\in\{1,2,3,4,6\}$, and we establish the
derivative\,$\leftrightarrow$\,negapolygamma duality at all orders: the
functional equation pairs $s=k+1$ with $s=-k$, the layer of the negative-order
polygamma $\psi^{(-(k+1))}$, through the harmonic-number bridge
$L'(k{+}1,\chi)/L(k{+}1,\chi)+\overline{L'(-k,\chi)/L(-k,\chi)}=\gamma+\ln(2\pi/q)-\h{k}$,
with a clean parity dichotomy. As a representative corollary,
$\gamma_1''(\tfrac14)=-256\pi^3\,\psi^{(-3)}(\tfrac14)+(\text{elementary})$.
\end{abstract}

\section{Introduction and notation}\label{sec:intro}

The generalized (Hurwitz) Stieltjes constants $\gamma_n(a)$ are the Laurent
coefficients of the Hurwitz zeta function at its pole,
\begin{equation}\label{eq:laurent}
  \zeta(s,a)=\frac{1}{s-1}+\sum_{n\ge0}\frac{(-1)^n}{n!}\,\gamma_n(a)\,(s-1)^n ,
\end{equation}
so $\gamma_0(a)=-\psi(a)$ and $\gamma_n(1)=\gamma_n$ is the ordinary $n$-th
Stieltjes constant. The word ``constant'' is a misnomer: each $\gamma_n$ is a
smooth function of $a$ on $(0,\infty)$, and the object of this note is the
\emph{parameter-derivative tower}
\[
  \gamma_n^{(k)}(a):=\frac{d^k}{da^k}\,\gamma_n(a),\qquad k=0,1,2,\dots
\]
Blagouchine~\cite{Blagouchine2015} gave the closed-form theory of $\gamma_1(p/q)$
at rational arguments; the present development is the analogue ``one (and two)
order(s) up,'' for the derivatives in the parameter, which had not been
tabulated.

Throughout, $\zeta^{(j)}(s_0,a)=\partial_s^j\zeta(s,a)\big|_{s=s_0}$ and
$L^{(j)}(s_0,\chi)$ similarly for a Dirichlet $L$-function; $\h{k}=\sum_{m=1}^k
1/m$, $\h{k}^{(2)}=\sum_{m=1}^k 1/m^2$ (with $\h0=\h0^{(2)}=0$); $\psi^{(k)}$ is
the $k$-th derivative of the digamma function and $\psi^{(-n)}$ the
negative-order polygamma (Wolfram's $\texttt{PolyGamma}[-n,\cdot]$); $\Glaisher$
is the Glaisher--Kinkelin constant, $\Cat$ Catalan's constant, $\beta$ the
Dirichlet beta function, and $\chim3,\chim4$ the odd quadratic characters of
conductors $3,4$.

\section{The uniform master identity}\label{sec:master}

\begin{theorem}[Harmonic-number master identity]\label{thm:master}
For all $k\ge1$ and $n\ge0$,
\begin{equation}\label{eq:masterN}
  \gamma_n^{(k)}(a)=(-1)^{n+k}\,k!\,n!\sum_{i+j=n} e_i(k)\,
   \frac{\zeta^{(j)}(k{+}1,a)}{j!},
\end{equation}
where $e_i(k)$ is the $i$-th elementary symmetric function of
$\{1,\tfrac12,\dots,\tfrac1k\}$; in particular $e_0=1$, $e_1=\h{k}$,
$e_2=\tfrac12(\h{k}^2-\h{k}^{(2)})$. The first two rows are
\begin{align}
  \gamma_1^{(k)}(a)&=(-1)^{k+1}k!\,\bigl[\zeta'(k{+}1,a)+\h{k}\,\zeta(k{+}1,a)\bigr],
   \label{eq:master1}\\
  \gamma_2^{(k)}(a)&=(-1)^{k}k!\,\bigl[\zeta''(k{+}1,a)+2\h{k}\,\zeta'(k{+}1,a)
   +(\h{k}^2-\h{k}^{(2)})\,\zeta(k{+}1,a)\bigr].\label{eq:master2}
\end{align}
\end{theorem}

\begin{proof}
From $\partial_a\zeta(s,a)=-s\,\zeta(s{+}1,a)$, iteration gives
$\partial_a^k\zeta(s,a)=(-1)^k (s)_k\,\zeta(s{+}k,a)$ with the rising factorial
$(s)_k=s(s{+}1)\cdots(s{+}k{-}1)$. For $k\ge1$ this annihilates the pole term in
\eqref{eq:laurent}, so
\[
  \sum_{n\ge0}\frac{(-1)^n}{n!}\,\gamma_n^{(k)}(a)\,(s-1)^n
   =(-1)^k(s)_k\,\zeta(s{+}k,a).
\]
Set $s=1+\varepsilon$. Then
$(s)_k=(1+\varepsilon)_k=\prod_{m=1}^k(m+\varepsilon)=k!\prod_{m=1}^k(1+\varepsilon/m)
 =k!\sum_{i\ge0}e_i(k)\,\varepsilon^i$,
the generating identity for the elementary symmetric functions of $\{1/m\}$,
while $\zeta(s{+}k,a)=\zeta(k{+}1{+}\varepsilon,a)=\sum_{j\ge0}\zeta^{(j)}(k{+}1,a)\,\varepsilon^j/j!$
is analytic at $\varepsilon=0$ because $k{+}1\ge2$. Comparing coefficients of
$\varepsilon^n$ yields \eqref{eq:masterN}; $n=1,2$ give
\eqref{eq:master1}--\eqref{eq:master2}.
\end{proof}

Theorem~\ref{thm:master} is the harmonic-number repackaging of Coffey's
all-orders Stirling-number ladder~\cite{Coffey2009}: the change of basis from
the unsigned Stirling numbers $\bigl|s(k{+}1,i{+}1)\bigr|$ to the symmetric
functions $e_i(k)$ turns the coefficients into $\h{k}$, $\h{k}^2-\h{k}^{(2)}$,
\dots. The specializations $\gamma_1'=\zeta'(2,a)+\zeta(2,a)$,
$\gamma_1''=-2\zeta'(3,a)-3\zeta(3,a)$,
$\gamma_2'=-[\zeta''(2,a)+2\zeta'(2,a)]$ and
$\gamma_2''=2\zeta''(3,a)+6\zeta'(3,a)+2\zeta(3,a)$ are the $s=2,3$ layers used
below. Identity~\eqref{eq:masterN} was verified against numerical $k$-th
parameter-derivatives of $\gamma_1,\gamma_2$ for $k=1,\dots,5$.

\section{The uniform rational-argument landscape}\label{sec:landscape}

Three structural facts hold uniformly in $k$; all follow from
$\sum_{p}\chi(p)\,\zeta(s,p/q)=q^{s}L(s,\chi)$ (and its multiplicative cousin)
differentiated in $s$, fed through Theorem~\ref{thm:master}.

\begin{proposition}[Character coordinate]\label{prop:char}
For a primitive Dirichlet character $\chi$ modulo $q$ and all $k\ge1$,
\begin{equation}\label{eq:char}
  \sum_{p=1}^{q-1}\chi(p)\,\gamma_1^{(k)}(p/q)
   =(-1)^{k+1}k!\,q^{k+1}\bigl[L'(k{+}1,\chi)+(\h{k}+\ln q)\,L(k{+}1,\chi)\bigr].
\end{equation}
\end{proposition}

\begin{proposition}[Distribution row]\label{prop:dist}
For $d\ge2$ and all $k\ge1$, with $\psi^{(k)}(a)=(-1)^{k+1}k!\,\zeta(k{+}1,a)$,
\begin{equation}\label{eq:dist}
  \sum_{j=0}^{d-1}\gamma_1^{(k)}\!\Bigl(\tfrac{a+j}{d}\Bigr)
   =d^{\,k+1}\gamma_1^{(k)}(a)+d^{\,k+1}\ln d\;\psi^{(k)}(a).
\end{equation}
\end{proposition}

\begin{proof}[Proof of \eqref{eq:char} and \eqref{eq:dist}]
Apply $\sum_p\chi(p)\,(\cdot)$ (resp. $\sum_{j<d}(\cdot)$) to
\eqref{eq:master1}. For \eqref{eq:char}, use
$\sum_p\chi(p)\zeta(s,p/q)=q^sL(s,\chi)$ and its $s$-derivative
$q^s\ln q\,L(s,\chi)+q^sL'(s,\chi)$ at $s=k{+}1$; collect. For \eqref{eq:dist},
use $\sum_{j<d}\zeta(s,(a{+}j)/d)=d^s\zeta(s,a)$ and its $s$-derivative, and
recognize $(-1)^{k+1}k!\,\zeta(k{+}1,a)=\psi^{(k)}(a)$.
\end{proof}

\begin{theorem}[$k$-independent rank]\label{thm:rank}
Fix a grid denominator $q$. The $\Q$-vector space of linear relations among
$\{\gamma_1^{(k)}(r/q):1\le r<q\}$ generated by the distribution
rows~\eqref{eq:dist} has, for every $k\ge1$, residual dimension
\[
  (q-1)-\operatorname{rank}=\varphi(q)-1 ,
\]
one residual direction per non-principal character mod $q$, both parities.
\end{theorem}

The rank was computed by exact rational row reduction; for $k=1,2,3$ and every
$q=3,\dots,30$ the residual vectors are bitwise identical (the collapsed
coordinate carries weight $d^{k+1}$, which changes the integer entries but not
the rank). Thus differentiation in the parameter destroys, at \emph{every}
order, the Malmsten--Blagouchine reflection halving $\varphi(q)/2-1$ of the
underived constants. Equation~\eqref{eq:char} was verified for
$\chim3,\chim4,\chi_5$ and $k=1,\dots,5$; \eqref{eq:dist} for
$k=1,\dots,5$, $d\in\{2,3,5\}$.

\section{Explicit tables for the second constant}\label{sec:tables}

We record the closed forms of $\gamma_2'=-[\zeta''(2,a)+2\zeta'(2,a)]$ (the
$s=2$ layer) and $\gamma_2''=2\zeta''(3,a)+6\zeta'(3,a)+2\zeta(3,a)$ (the $s=3$
layer) at the elementary denominators. The new constants entering are the
\emph{second} $s$-derivatives $\zeta''(2),\zeta''(3)$, $\beta''(s)=L''(s,\chim4)$
and $L''(s,\chim3)$; the values $L(3,\chim4)=\pi^3/32$ and
$L(3,\chim3)=4\pi^3/(81\sqrt3)$ are elementary. As at the first-derivative
layers, the odd-character block flips sign as a single unit under $p\mapsto q-p$
(including any elementary $\pi$-power value-part it carries). All entries are
exact and were verified to $<10^{-70}$.

\paragraph{Second constant, first parameter-derivative.}
\begin{align*}
\gamma_2'(1)   &= -\zeta''(2)-2\zeta'(2),\\
\gamma_2'(\tfrac12) &= -3\zeta''(2)-(6+8\ln2)\zeta'(2)-\tfrac43\pi^2\ln2-\tfrac23\pi^2\ln^2 2,\\
\gamma_2'(\tfrac13) &= -4\zeta''(2)-(8+9\ln3)\zeta'(2)-\tfrac32\pi^2\ln3-\tfrac34\pi^2\ln^2 3\\
 &\quad -\Bigl[\tfrac92 L''(2,\chim3)+9(1{+}\ln3)L'(2,\chim3)+\bigl(9\ln3+\tfrac92\ln^2 3\bigr)L(2,\chim3)\Bigr],\\
\gamma_2'(\tfrac14) &= -6\zeta''(2)-(12+28\ln2)\zeta'(2)-\tfrac{14}{3}\pi^2\ln2-5\pi^2\ln^2 2\\
 &\quad -\Bigl[8\beta''(2)+16(1{+}2\ln2)\beta'(2)+\bigl(32\ln2+32\ln^2 2\bigr)\Cat\Bigr],
\end{align*}
with $\gamma_2'(2/3),\gamma_2'(3/4)$ obtained by negating the bracketed
odd-character block, and the $q=6$ entries given in the source report.

\paragraph{Second constant, second parameter-derivative.}
\begin{align*}
\gamma_2''(1)   &= 2\zeta''(3)+6\zeta'(3)+2\zeta(3),\\
\gamma_2''(\tfrac12) &= 14\zeta''(3)+(42+32\ln2)\zeta'(3)+(14+48\ln2+16\ln^2 2)\zeta(3),\\
\gamma_2''(\tfrac13) &= 26\zeta''(3)+(78+54\ln3)\zeta'(3)+(26+81\ln3+27\ln^2 3)\zeta(3)\\
 &\quad +\Bigl[27L''(3,\chim3)+(81+54\ln3)L'(3,\chim3)+(27+81\ln3+27\ln^2 3)L(3,\chim3)\Bigr],\\
\gamma_2''(\tfrac14) &= 56\zeta''(3)+(168+240\ln2)\zeta'(3)+(56+360\ln2+248\ln^2 2)\zeta(3)\\
 &\quad +\Bigl[64\beta''(3)+(192+256\ln2)\beta'(3)+(2+12\ln2+8\ln^2 2)\pi^3\Bigr],
\end{align*}
again with $\gamma_2''(2/3),\gamma_2''(3/4)$ the odd-block negations. The
reducible fractions collapse correctly ($\tfrac24=\tfrac36=\tfrac12$,
$\tfrac26=\tfrac13$), an internal consistency check.

\section{The derivative\,$\leftrightarrow$\,negapolygamma duality}\label{sec:duality}

The functional equation pairs the layer $s=k+1$ of $\gamma_1^{(k)}$ with the
layer $s=-k$, which is exactly where the negative-order polygamma
$\psi^{(-(k+1))}$ lives (its ladder uses $\zeta'(-k,\cdot)$). The seam is a
single uniform identity.

\begin{theorem}[Harmonic-number bridge law]\label{thm:bridge}
Let $\chi$ be a primitive character mod $q$ for which $L(-k,\chi)\neq0$. Then
\begin{equation}\label{eq:bridge}
  \frac{L'(k{+}1,\chi)}{L(k{+}1,\chi)}
   +\overline{\left(\frac{L'(-k,\chi)}{L(-k,\chi)}\right)}
   =\gamma+\ln\frac{2\pi}{q}-\h{k}.
\end{equation}
The same $\h{k}$ that sets the master coefficient in
Theorem~\ref{thm:master} sets the bridge constant; the case $k=1$ recovers the
generalized-Glaisher constant $C_q=\gamma+\ln(2\pi/q)-1$.
\end{theorem}

\begin{remark}[Parity dichotomy]\label{rem:parity}
Since $L(-k,\chi)=0$ if and only if $\chi(-1)=(-1)^k$, the
direct form \eqref{eq:bridge} applies precisely when $\chi(-1)\neq(-1)^k$;
otherwise the $s=-k$ side has a trivial zero and \eqref{eq:bridge} holds in its
first-nonzero-order variant (a ratio of leading nonzero derivatives). This is
what makes the \emph{value-level} corollary parity-selective.
\end{remark}

At $z=\tfrac14$ the governing character is $\chim4$ (odd). For $k$ even there is
no trivial zero, and $\psi^{(-(k+1))}(\tfrac14)$ carries the genuine derivative
atom $\beta'(k{+}1)=L'(k{+}1,\chim4)$, so $\gamma_1^{(k)}(\tfrac14)$ is an exact
finite combination of $\psi^{(-(k+1))}(\tfrac14)$ and elementary constants. The
representative instance $k=2$ is, after eliminating $\beta'(3)$ entirely,
\begin{equation}\label{eq:k2duality}
\begin{aligned}
  \gamma_1''(\tfrac14)=\;&-256\,\pi^3\,\psi^{(-3)}(\tfrac14)
   -56\,\zeta'(3)-84\,\zeta(3)-120\ln2\,\zeta(3)+35\pi\,\zeta(3)\\
   &-3\pi^3-4\ln2\,\pi^3+64\ln\Glaisher\,\pi^3-2\gamma\,\pi^3+2\ln(2\pi)\,\pi^3 ,
\end{aligned}
\end{equation}
verified to $<10^{-118}$. (The irrational coefficient $-256\pi^3$ on
$\psi^{(-3)}(\tfrac14)$ is why a rational integer-relation search for this
identity must take $\pi^3\psi^{(-3)}$, not $\psi^{(-3)}$, as the search atom.)
For $k$ odd the value-level duality \emph{degenerates}: by
Remark~\ref{rem:parity}, $\psi^{(-(k+1))}(\tfrac14)$ then carries only the lower
\emph{value} atom (for $k=1$, $\beta'(-1)=2\Cat/\pi$, elementary; for $k=3$,
$\beta'(-3)\sim\beta(4)/\pi^3$), while the genuine derivative atom
$\beta'(k{+}1)$ of $\gamma_1^{(k)}(\tfrac14)$ is tied instead to the second
derivative $\beta''(-k)$, absent from $\psi^{(-(k+1))}$. Numerically the
coefficient of $\beta'(k{+}1)$ in $\psi^{(-(k+1))}(\tfrac14)$ is zero to $\sim
10^{-163}$ for $k=1,3$. The clean odd-$k$ value-level duality therefore lives at
even-character points (e.g.\ $q=5$), where the parity is reversed.

\section{The new atoms are genuine}\label{sec:atoms}

The $\gamma_2$ tables of \S\ref{sec:tables} introduce six second-$s$-derivative
constants:
$\zeta''(2),\zeta''(3),\beta''(2),\beta''(3),L''(2,\chim3),L''(3,\chim3)$. A
PSLQ sweep at $200$ digits with planted controls (the elementary values
$\beta(3)=\pi^3/32$, $L(3,\chim3)=4\pi^3/(81\sqrt3)$, $\zeta(2)=\pi^2/6$ are
recovered; a random constant is rejected) finds \emph{no} reduction of any of
the six over a basket of the lower value and first-derivative atoms with
$\gamma$/log products, and no relation among them: each is a genuine,
mutually independent atom, the natural continuation of the $\beta'(2),\beta'(3)$
findings at the first-derivative layers. This was reproduced by an independent
agent with its own controls.

\section*{Verification record}

Every identity above was checked numerically in the
\emph{Smithereens} repository: the master identity, character coordinate,
distribution row, rank theorem, $\gamma_2$ tables, bridge law and the
duality~\eqref{eq:k2duality} by the self-test battery in
\texttt{src/PolyLog/tools/stieltjes\_general\_k.py} (with
\texttt{scratch/g1k\_grid\_rank.py} for the exact rank and
\texttt{stieltjes\_s2deriv\_hunts.py} for \S\ref{sec:atoms}); the $q=4$ table
entries and the master identities additionally in an independent
computer-algebra system. An independent multi-agent workflow re-derived the
$\gamma_2''$ table (all entries to $<10^{-128}$), re-verified the general-$k$
claims at the independent point $a=\sqrt2-1$ for $k=1,\dots,5$, confirmed the
bridge law in both parities ($\le 2.7\times10^{-161}$), and reproduced the atom
sweep.

\begin{thebibliography}{9}
\bibitem{Blagouchine2015} I.~V.~Blagouchine,
  \emph{A theorem for the closed-form evaluation of the first generalized
  Stieltjes constant at rational arguments and some related summations},
  J.\ Number Theory \textbf{148} (2015), 537--592.
\bibitem{Coffey2009} M.~W.~Coffey,
  \emph{Series representations for the Stieltjes constants},
  arXiv:0905.1111; Rocky Mountain J.\ Math.\ \textbf{44} (2014), 443--477.
\bibitem{Connon2019} D.~F.~Connon,
  \emph{On some new formulae involving the Stieltjes constants} (2019).
\bibitem{MillerAdamchik1998} J.~Miller and V.~S.~Adamchik,
  \emph{Derivatives of the Hurwitz zeta function for rational arguments},
  J.\ Comput.\ Appl.\ Math.\ \textbf{100} (1998), 201--206.
\bibitem{BaileyBorwein2016} D.~H.~Bailey and J.~M.~Borwein,
  \emph{Computation and structure of character polylogarithms with
  applications to character Mordell--Tornheim--Witten sums},
  Math.\ Comp.\ \textbf{85} (2016), 295--324.
\end{thebibliography}

\end{document}
